\documentclass[10pt,a4paper]{amsart}
\usepackage[margin=19mm]{geometry}
\usepackage{amsmath,amssymb,mathtools}
\usepackage[scr=rsfs,cal=euler]{mathalfa}
\usepackage{xfrac}
\usepackage[unicode,hidelinks,bookmarks=false]{hyperref}
\newcommand{\ie}{i.e.\ }
\newcommand{\cf}{cf.\ }
\newcommand{\resp}{respectively\ }
\newcommand{\Subsection}{\subsection}
\providecommand{\nobreakdash}{\nobreak\hbox{-}}
\setlength{\emergencystretch}{2em}
\multlinegap0pt
\usepackage{amssymb}
\usepackage{bm}
\usepackage{float}
\usepackage{xcolor}
\newtheorem{theorem}{Theorem}
\title{Perfect numbers, Wieferich primes and the solutions of \texorpdfstring{$\displaystyle\binom{2n}{n}\equiv 2^n \bmod n$} {binomial (2n,n) congruent to 2^n modulo n}}
\author{Gabriel Guedes, Ricardo Machado}
\date{Source: arXiv:2609.03091v1 (CC BY 4.0)}
\begin{document}
\maketitle

\noindent\emph{Source and licence.} Perfect numbers, Wieferich primes and the solutions of \texorpdfstring{$\displaystyle\binom{2n}{n}\equiv 2^n \bmod n$} {binomial (2n,n) congruent to 2^n modulo n}. Version arXiv:2609.03091v1 (CC BY 4.0), distributed by arXiv under the Creative Commons Attribution 4.0 International licence, http://creativecommons.org/licenses/by/4.0/, as recorded in the arXiv licence metadata for this exact version and shown on the abstract page https://arxiv.org/abs/2609.03091v1. Full licence notice: \texttt{licenses/arxiv-2609.03091-CC-BY-4.0.txt}.\par\smallskip

\noindent\textit{Editorial scope.} Counted unit: the article's theorem teo-32 (source0.tex lines 311--314). The article's complete original proof of it is delivered with it (source0.tex lines 316--339, one \texttt{proof} environment of 24 lines). No mathematics is written, repaired or re-derived: the statement and the proof are the article's own lines, in the article's own notation, and no retained line is substituted or re-flowed.  Retained uncounted as the source-local context the proof needs: lines 143--145.  Nothing was omitted from any retained span, so the package carries no omission record.  No retained line contains a citation, so the wrapper carries no bibliography and no bibliography entry.  No retained line contains a figure environment, an image-inclusion command or drawing code, so no image is delivered and the package declares no asset dependency.  Editorial formatting only, in addition: an AMS \texttt{amsart} preamble that restates the article's own declarations and macros used by the retained lines, namely the environments \texttt{theorem} on separate counters, together with the standard packages those lines need.  The retained environments print in this document as Theorem 1 in order of appearance, whereas the article prints the same items with the numbering of its own full document.  The complete native source of the article as distributed in the arXiv e-print for this exact version is bundled unchanged as \texttt{sources/209263/source0.tex}.
\begin{equation} \label{eq-001}
\binom{2n}{n} \equiv 2^n \bmod n
\end{equation}

\begin{theorem}\label{teo-32}
Let $n=p^2$ with $p$ a prime number. We have that
$p$ is a Wieferich prime if and only if  $n$ is a solution of Equation~\eqref{eq-001}.
\end{theorem}

\begin{proof}
Supposing that $p$ is a Wieferich prime we are going to prove that $n$ is a solution of Equation \eqref{eq-001}, that is,
\begin{equation}\label{eq8}
\binom{2p^2}{p^2}\equiv2^{p^2} \bmod {p^2}.
\end{equation}
Using Lucas's Theorem  we have that left-hand side of Equation \eqref{eq8} is
$$\binom{2p^2}{p^2}\equiv\binom{2}{1}\equiv 2 \bmod {p^2}.$$ In the right-hand side of \eqref{eq8} note that
$2^{p^2}=2^{p(p-1)}\cdot2^p$,
by Euler's Theorem 
$2^{p(p-1)}\equiv 1 \bmod {p^2}$
then 
$$2^{p^2}\equiv2^{p(p-1)}\cdot2^p\equiv2^p\equiv2^p -2 +2\equiv 2\cdot(2^{p-1}-1) +2 \bmod {p^2}.$$
Now apply the hypothesis of $p$ being a Wieferich prime then
$$2\cdot(2^{p-1}-1) +2\equiv 2 \bmod {p^2.}$$
Thus, we show that Equation \eqref{eq8} is valid, since both sides are congruent to $2$ modulo $p^2$.


For the converse in the same lines we have that
$2^p\equiv 2 \bmod p^2$
simplifying by $2$ we obtain
$2^{p-1}\equiv 1 \bmod {p^2}$
which % TECH EDITOR: wich > which
is the definition of a Wieferich prime.
\end{proof}

\end{document}
